\section{환경 생성과 능동 실패 탐색}

\subsection*{장면 구조화와 후보 생성}
환경은 객체(목표물, 장애물, 목적지, 사람)의 위치·크기·회전·물성으로 이루어진 파라미터
벡터로 정의된다. 파라미터가 바뀌면 3D 장면 그래프, 이동 지도, MuJoCo의 XML 형식(MJCF)~\cite{mjcf}으로 된 물리 시뮬레이션 모델을 함께 생성하고
서로 일치하는지 검사한다. 로봇팔 환경은 그릇·접시 위치와 장애물 상자의 위치·크기·높이·회전
등 \textbf{10개 축}, 휴머노이드 환경은 상자·바닥 물성, 통로 폭, 장애물 2개의 형상, 사람의 이동을
포함한 \textbf{21개 축}을 제공한다(부록 A). 생성된 후보는 객체 간 관통, 작업공간 범위, 로봇
몸체와의 겹침을 로봇 실행 전에 검사하여 의미 없는 테스트를 제거한다.

\subsection*{행동 피드백 기반 탐색 절차}
\begin{enumerate}
\item 이전 실행의 목표 결과, 실제 행동, 로봇 기능 호출 결과(경로 계획, 밀기 등), 접촉·이동량 측정, 종료 상태를 요약한다.
\item 탐색용 LLM이 관측 사실을 근거로 실패 가설을 세우고, 허용 범위 안의 환경 변경값을 정해진 형식으로 제안한다.
\item 후보 검증기가 값의 범위, 근거로 든 이전 실행의 실재 여부, 변경 축 수, 예산을 검사한다. 범위 밖 값은 실행 전에 거부한다.
\item 새 환경에서 로봇을 실행하고 독립된 평가기가 판정한다. 예상과 실제의 차이를 다음 제안에 반영한다.
\end{enumerate}
성공 기록도 함께 제공하여, 이미 알려진 실패만 반복하지 않고 성공 경계 근처를 탐색하도록 한다.
Random은 같은 범위에서 축과 값을 균등하게 무작위 추출하며 결과를 보고 재추출하지 않는다.

\subsection*{실제 탐색 연결 예}
표~\ref{tab:link}는 두 실험에서 이전 실행의 관측이 LLM의 가설과 다음 환경 제안을 거쳐 실제 실패로 이어진 기록이다.
\begin{table*}[!t]
\centering\small
\caption{관측에서 다음 시험으로 이어진 실제 기록}\label{tab:link}
\begin{tabularx}{\textwidth}{@{}p{22mm}YY@{}}
\toprule
단계 & 로봇팔 (시드 661) & 휴머노이드 \\
\midrule
이전 실행 & 초기 장면 FAIL: 로봇–장애물 접촉 13.3초, 그릇 미이동 & 기준 장면 PASS: 연속 이동으로 목표 도달 \\
LLM 가설 & 이송 경로 중앙의 더 크고 높은 상자가 접근을 막을 것 & 상자를 목표 쪽으로 옮기면 목표 영역이 점유될 것 \\
제안 & 장애물 폭·깊이 6\,cm, 높이 17\,cm & 상자를 통로 가장자리에서 목표 위치로 이동 (1개 축 변경) \\
결과 & FAIL, 접촉 11.0초, 그릇 미이동 & FAIL, 목표 위치 점유 검출 \\
\bottomrule
\end{tabularx}
\end{table*}
